\documentclass[11pt, a4paper]{article}

\usepackage[a4paper, margin=2.5cm]{geometry}
\usepackage{fontspec}
\usepackage[russian, bidi=basic, provide=*]{babel}
\babelprovide[import, onchar=ids fonts]{russian}
\babelprovide[import, onchar=ids fonts]{english}
\babelfont{rm}{Noto Sans}
\babelfont{sf}{Noto Sans}
\babelfont{tt}{Noto Sans Mono}

\usepackage{enumitem}
\setlist[itemize]{label=-}
\usepackage{amsmath}
\usepackage{booktabs}
\usepackage{tabularx}
\usepackage{hyperref}
\usepackage{listings}
\usepackage{xcolor}
\usepackage[htt]{hyphenat}

\definecolor{codegray}{rgb}{0.5,0.5,0.5}
\definecolor{backcolor}{rgb}{0.97,0.97,0.97}

\lstset{
    backgroundcolor=\color{backcolor},
    basicstyle=\ttfamily\scriptsize,
    breaklines=true,
    breakatwhitespace=false,
    breakautoindent=true,
    postbreak=\mbox{\textcolor{codegray}{\(\hookrightarrow\)\ }},
    columns=fullflexible,
    keepspaces=true,
    frame=single,
    rulecolor=\color{black!20},
    numbers=left,
    numbersep=6pt,
    numberstyle=\tiny\color{codegray},
    showstringspaces=false,
    tabsize=2,
    xleftmargin=15pt,
    framexleftmargin=15pt
}

\hypersetup{
    colorlinks=true,
    linkcolor=black,
    urlcolor=blue
}

\begin{document}

\begin{titlepage}
    \centering
    {\small МИНИСТЕРСТВО НАУКИ И ВЫСШЕГО ОБРАЗОВАНИЯ РОССИЙСКОЙ ФЕДЕРАЦИИ\par}
    \vspace{0.3cm}
    {\small Федеральное государственное автономное образовательное учреждение высшего образования\par}
    {\small «Московский авиационный институт (национальный исследовательский университет)»\par}
    
    \vspace{4cm}
    {\Large \textbf{ОТЧЁТ ПО ЛАБОРАТОРНОЙ РАБОТЕ №2}\par}
    
    \vspace{1.2cm}
    {\large \textbf{Дисциплина:}\par}
    {\large Операционные системы и архитектура компьютера\par}
    
    \vspace{0.8cm}
    {\large \textbf{Тема:}\par}
    {\large Управление процессами в ОС. Обеспечение обмена данными между процессами посредством каналов\par}
    
    \vspace{3.5cm}
    \begin{flushright}
        \begin{minipage}{0.48\textwidth}
            \textbf{Выполнил:}\\
            Студент группы М8О-203БВ-25\\
            К. П. Смаге\\[0.8cm]
            \textbf{Преподаватель:}\\
            Е. С. Миронов\\
        \end{minipage}
    \end{flushright}
    
    \vfill
    {\large Москва, 2026 г.\par}
\end{titlepage}

\section{Цель работы}
Приобретение практических навыков в:
\begin{itemize}
    \item управлении процессами в операционных системах семейств Linux (POSIX) и Windows (Win32 API);
    \item обеспечении межпроцессного взаимодействия (IPC) посредством неименованных каналов (\texttt{pipe}).
\end{itemize}

\section{Краткое описание архитектуры решения}
Разработанное решение реализует требования варианта 20. Программный комплекс состоит из двух отдельных исполняемых файлов: управляющей родительской программы (\texttt{parent}) и рабочей дочерней программы (\texttt{child}).

Алгоритм функционирования системы включает следующие этапы:
\begin{enumerate}
    \item Родительский процесс считывает со стандартного потока ввода две строки, содержащие пути к выходным файлам для первого и второго дочерних процессов.
    \item Родительский процесс инициализирует два экземпляра класса \texttt{Process} (\texttt{child1} и \texttt{child2}). Для каждого дочернего процесса создается однонаправленный неименованный канал (\texttt{pipe1} и \texttt{pipe2}).
    \item Стандартный поток ввода (\texttt{stdin}) каждого дочернего процесса перенаправляется на читающий конец соответствующего канала, а стандартный поток вывода (\texttt{stdout}) перенаправляется в целевой выходной файл на диске. Стандартный поток ошибок (\texttt{stderr}) остается привязанным к терминалу для вывода диагностических сообщений и идентификаторов процессов (PID).
    \item Родительский процесс построчно считывает пользовательский ввод и осуществляет фильтрацию согласно условию варианта 20:
    \begin{itemize}
        \item если длина строки превышает 10 символов ($\text{length} > 10$), строка пересылается через \texttt{pipe2} процессу \texttt{child2};
        \item если длина строки не превышает 10 символов ($\text{length} \le 10$), строка пересылается через \texttt{pipe1} процессу \texttt{child1}.
    \end{itemize}
    \item Каждый дочерний процесс в цикле считывает поступающие строки из своего стандартного ввода, инвертирует порядок символов в строке и записывает результат в стандартный поток вывода (целевой файл).
    \item При получении сигнала конца ввода (\texttt{EOF}) родительский процесс закрывает пишущие дескрипторы обоих каналов, сигнализируя дочерним процессам о завершении потока данных, и ожидает их штатного завершения, предотвращая образование процессов-зомби.
\end{enumerate}

Для достижения кроссплатформенности применен паттерн изоляции платформенного слоя: бизнес-логика в \texttt{src/parent.cpp} и \texttt{src/child.cpp} не содержит прямых системных вызовов ОС и взаимодействует с ресурсами ядра исключительно через абстракцию \texttt{Process} (\texttt{include/process.hpp}), реализующую идиому RAII. Выбор конкретной реализации (\texttt{src/posix/process\_posix.cpp} для Linux или \texttt{src/win/process\_win.cpp} для Windows) выполняется системой сборки CMake на этапе конфигурации.

В таблице 1 приведено сопоставление системных примитивов, используемых в платформенных модулях для Linux и Windows.

\begin{table}[h!]
\centering
\small
\begin{tabularx}{\textwidth}{l X X}
\toprule
\textbf{Операция ОС} & \textbf{Linux (POSIX API)} & \textbf{Windows (Win32 API)} \\
\midrule
Создание канала & \texttt{pipe(pipefd)} & \texttt{CreatePipe(\&pipefd[0], \&pipefd[1], \&sa, 0)} \\
Открытие файла вывода & \texttt{open(path, O\_WRONLY | O\_CREAT | O\_TRUNC, 0644)} & \texttt{CreateFileA(path, GENERIC\_WRITE, ..., CREATE\_ALWAYS, ...)} \\
Перенаправление потоков & \texttt{dup2(fd, STDIN\_FILENO / STDOUT\_FILENO)} & Поля \texttt{hStdInput}, \texttt{hStdOutput} в \texttt{STARTUPINFOA} \\
Изоляция конца записи & \texttt{close(pipefd[1])} в дочернем процессе & \texttt{SetHandleInformation(pipefd[1], HANDLE\_FLAG\_INHERIT, 0)} \\
Порождение процесса & \texttt{fork()} + \texttt{execl(path, path, nullptr)} & \texttt{CreateProcessA(nullptr, cmd.data(), ..., \&si, \&pi)} \\
Запись в канал & \texttt{write(write\_fd, buf, size)} & \texttt{WriteFile(write\_fd, buf, size, \&written, nullptr)} \\
Ожидание завершения & \texttt{waitpid(pid, \&status, 0)} & \texttt{WaitForSingleObject(pid, INFINITE)} \\
Освобождение ресурсов & \texttt{close(fd)} & \texttt{CloseHandle(handle)} \\
\bottomrule
\end{tabularx}
\caption{Сопоставление системных вызовов в реализациях для Linux и Windows}
\end{table}

\section{Листинг ключевых фрагментов кода}

\subsection{Кроссплатформенный интерфейс управления процессом}
Заголовочный файл \texttt{include/process.hpp} инкапсулирует платформенные типы дескрипторов (\texttt{int} и \texttt{pid\_t} в POSIX, \texttt{HANDLE} в Win32 API) и предоставляет единый RAII-интерфейс для запуска дочернего процесса, записи строк в канал и синхронизации завершения.

\begin{lstlisting}
#ifdef _WIN32
    #include <windows.h>
    using pipe_handle_t = HANDLE;
    using proc_handle_t = HANDLE;
#else
    #include <unistd.h>
    #include <sys/types.h>
    using pipe_handle_t = int;
    using proc_handle_t = pid_t;
#endif

class Process {
public:
    Process(
        const std::string& binary_path, 
        const std::string& output_file_path
    );

    ~Process();

    Process(const Process&) = delete;
    Process& operator=(const Process&) = delete;

    void send_line(const std::string& line);
    void close_input();
    void wait();

    unsigned long get_pid() const;

private:
    pipe_handle_t write_fd;
    proc_handle_t pid;
    bool is_closed;
};
\end{lstlisting}
\begin{center}
    \small Листинг 1 — Объявление класса \texttt{Process} (\texttt{include/process.hpp})
\end{center}

\subsection{Маршрутизация данных в родительском процессе}
Функция \texttt{parent()} в файле \texttt{src/parent.cpp} отвечает за чтение имен файлов, создание двух дочерних процессов, распределение входных строк по критерию длины ($\le 10$ или $> 10$ символов) и корректное закрытие каналов.

\begin{lstlisting}
void parent() {
    std::cout << "Enter output file for Child1: ";
    std::string file1{};
    if (!std::getline(std::cin, file1) || file1.empty()) {
        throw std::runtime_error("Failed to read filename");
    }

    std::cout << "Enter output file for Child2: ";
    std::string file2{};
    if (!std::getline(std::cin, file2) || file2.empty()) {
        throw std::runtime_error("Failed to read filename");
    }

    Process child1(CHILD_PATH, file1);
    Process child2(CHILD_PATH, file2);

    std::string line{};
    
    std::cout << "Enter lines for processing (Ctrl+D / Ctrl+C to finish): \n";
    while (std::getline(std::cin, line)) {
        if (line.length() > 10) {
            child2.send_line(line);
        } else {
            child1.send_line(line);
        }
    }

    child1.close_input();
    child2.close_input();

    child1.wait();
    child2.wait();
}

\end{lstlisting}
\begin{center}
    \small Листинг 2 — Логика диспетчеризации строк в родительском процессе (\texttt{src/parent.cpp})
\end{center}

\subsection{Реализация IPC и порождения процесса в ОС Linux (POSIX)}
В файле \texttt{src/posix/process\_posix.cpp} создание процесса выполняется связкой вызовов \texttt{pipe()}, \texttt{fork()}, \texttt{open()}, \texttt{dup2()} и \texttt{execl()}. В каждом системном вызове проверяется код возврата с выводом ошибки через \texttt{perror}.

\begin{lstlisting}
Process::Process(
        const std::string& binary_path, 
        const std::string& output_file_path
    ) {
        int file_fd = open(
            output_file_path.c_str(), 
            O_WRONLY | O_CREAT | O_TRUNC, 
            0644
        );

        if (file_fd == -1) {
            perror("open");
            throw std::runtime_error("Failed to open output file: " + output_file_path);
        }
        
        int pipefd[2]{};
        
        if (pipe(pipefd) == -1) {
            perror("pipe");
            close(file_fd);
            
            throw std::runtime_error("Failed to create pipe");
        }

        pid_t fork_rc = fork();
        
        if (fork_rc == -1) {
            // failed
            perror("fork");

            close(pipefd[0]);
            close(pipefd[1]);

            throw std::runtime_error("Failed to fork");
        } else if (fork_rc == 0) {
            // child
            if (dup2(pipefd[0], STDIN_FILENO) == -1) {
                perror("dup2 stdin");
                _exit(EXIT_FAILURE);
            }

            if (dup2(file_fd, STDOUT_FILENO) == -1) {
                perror("dup2 stdout");
                _exit(EXIT_FAILURE);
            }

            close(file_fd);
            close(pipefd[0]);
            close(pipefd[1]);

            execl(
                binary_path.c_str(), 
                binary_path.c_str(), 
                static_cast<char*>(nullptr)
            );

            perror("execl");
            _exit(127);
        } else {
            // parent
            this->pid = fork_rc;
            this->write_fd = pipefd[1];
            this->is_closed = false;

            close(pipefd[0]);
            close(file_fd);
        }
}

void Process::send_line(const std::string& line) {
    if (is_closed || write_fd == -1) {
        throw std::runtime_error("Attempted to write to a closed or invalid pipe");
    }

    std::string f_line = line;
    
    if (f_line.empty() || f_line.back() != '\n') {
        f_line += '\n';
    }

    ssize_t bytes_written = write(write_fd, f_line.data(), f_line.size());
    
    if (bytes_written == -1) {
        perror("write");
        throw std::runtime_error("Failed to write data to pipe");
    }
}
\end{lstlisting}
\begin{center}
    \small Листинг 3 — Порождение процесса и перенаправление дескрипторов в POSIX (\texttt{src/posix/process\_posix.cpp})
\end{center}

\subsection{Реализация IPC и порождения процесса в ОС Windows (Win32 API)}
В файле \texttt{src/win/process\_win.cpp} настройка наследуемости дескрипторов (\texttt{SECURITY\_ATTRIBUTES}), изоляция пишущего конца трубы (\texttt{SetHandleInformation}) и привязка стандартных потоков (\texttt{STARTUPINFOA}) выполняются в родительском процессе перед вызовом \texttt{CreateProcessA}.

\begin{lstlisting}
Process::Process(
    const std::string& binary_path, 
    const std::string& output_file_path
) {
    SECURITY_ATTRIBUTES sa{};
    sa.nLength = sizeof(SECURITY_ATTRIBUTES);
    sa.bInheritHandle = TRUE;
    sa.lpSecurityDescriptor = nullptr;

    HANDLE file_fd = CreateFileA(
        output_file_path.c_str(),
        GENERIC_WRITE,
        FILE_SHARE_READ,
        &sa,
        CREATE_ALWAYS,
        FILE_ATTRIBUTE_NORMAL,
        nullptr
    );
    if (file_fd == INVALID_HANDLE_VALUE) {
        std::cerr << "CreateFileA error: " << GetLastError() << '\n';
        throw std::runtime_error("Failed to open output file: " + output_file_path);
    }

    HANDLE pipefd[2]{INVALID_HANDLE_VALUE, INVALID_HANDLE_VALUE};
    if (!CreatePipe(&pipefd[0], &pipefd[1], &sa, 0)) {
        std::cerr << "CreatePipe error: " << GetLastError() << '\n';
        CloseHandle(file_fd);
        throw std::runtime_error("Failed to create pipe");
    }

    if (!SetHandleInformation(pipefd[1], HANDLE_FLAG_INHERIT, 0)) {
        std::cerr << "SetHandleInformation error: " << GetLastError() << '\n';
        CloseHandle(file_fd);
        CloseHandle(pipefd[0]);
        CloseHandle(pipefd[1]);
        throw std::runtime_error("Failed to set handle information");
    }

    STARTUPINFOA si{};
    si.cb = sizeof(STARTUPINFOA);
    si.dwFlags = STARTF_USESTDHANDLES;
    si.hStdInput = pipefd[0];
    si.hStdOutput = file_fd;
    si.hStdError = GetStdHandle(STD_ERROR_HANDLE);

    PROCESS_INFORMATION pi{};
    std::string cmd = binary_path;

    BOOL create_rc = CreateProcessA(
        nullptr,
        cmd.data(),
        nullptr,
        nullptr,
        TRUE,
        0,
        nullptr,
        nullptr,
        &si,
        &pi
    );

    if (!create_rc) {
        std::cerr << "CreateProcessA error: " << GetLastError() << '\n';
        CloseHandle(file_fd);
        CloseHandle(pipefd[0]);
        CloseHandle(pipefd[1]);
        throw std::runtime_error("Failed to create process");
    }

    this->pid = pi.hProcess;
    this->write_fd = pipefd[1];
    this->is_closed = false;

    CloseHandle(pi.hThread);
    CloseHandle(pipefd[0]);
    CloseHandle(file_fd);
}

void Process::send_line(const std::string& line) {
    if (is_closed || write_fd == INVALID_HANDLE_VALUE) {
        throw std::runtime_error("Attempted to write to a closed or invalid pipe");
    }

    std::string f_line = line;
    if (f_line.empty() || f_line.back() != '\n') {
        f_line += '\n';
    }

    DWORD bytes_written = 0;
    BOOL write_rc = WriteFile(
        write_fd,
        f_line.data(),
        static_cast<DWORD>(f_line.size()),
        &bytes_written,
        nullptr
    );

    if (!write_rc) {
        std::cerr << "WriteFile error: " << GetLastError() << '\n';
        throw std::runtime_error("Failed to write data to pipe");
    }
}

\end{lstlisting}
\begin{center}
    \small Листинг 4 — Порождение процесса и настройка каналов в Windows (\texttt{src/win/process\_win.cpp})
\end{center}

\subsection{Обработка строк в дочернем процессе}
Дочерний процесс (\texttt{src/child.cpp}) считывает строки из \texttt{std::cin}, инвертирует их функцией \texttt{reverse\_string} и записывает результат в перенаправленный поток \texttt{std::cout}.

\begin{lstlisting}
void reverse_string(std::string& line) {
    std::reverse(line.begin(), line.end());
}

void child() {
    std::string line{};

    while (std::getline(std::cin, line)) {
        reverse_string(line);
        std::cout << line << std::endl;
    }
}
\end{lstlisting}
\begin{center}
    \small Листинг 5 — Инвертирование строк и вывод телеметрии в \texttt{src/child.cpp}
\end{center}

\section{Результаты работы}

\subsection{Консольный вывод и логи выполнения}
Ниже приведен фактический протокол работы программы в терминале с отображением идентификаторов процессов (\texttt{PID}), подтверждающий параллельное существование родительского и двух дочерних процессов, корректную маршрутизацию строк по длине и штатное завершение по сигналу \texttt{EOF}:

\begin{lstlisting}
$ ./parent
[Parent] PID: 463
Enter output file for Child1: short_lines.txt
Enter output file for Child2: long_lines.txt
[Parent] Created Child1 (PID: 479)
[Parent] Created Child2 (PID: 480)
Enter lines for processing (Ctrl+D / Ctrl+C to finish): 
12345 
1234567
12345678910
123456789101112
12
123456789876543210
[Parent] Closed Child1
[Parent] Closed Child2

$ cat short_lines.txt 
54321
7654321
21

$ cat long_lines.txt 
01987654321
211101987654321
012345678987654321
\end{lstlisting}

\section{Выводы}
В результате выполнения лабораторной работы реализован кроссплатформенный программный комплекс межпроцессного взаимодействия, осуществляющий потоковую фильтрацию и инвертирование строк с использованием каналов. В ходе разработки были учтены архитектурные различия между моделями управления процессами в рассматриваемых ОС. 

В процессе работы мною были изучены и отработаны навыки перенаправления потоков ввода/вывода между процессами и практика безопасного управления ресурсами.

\end{document}
